%%%%% O1 -> ch4_satr.tex (inserted before \section{Summary})
\section{Specialisation to AppWorld}
\label{sec:satr-appworld}

On AppWorld the session is one task episode, and the catalogue is the 457 APIs of the nine applications.
SATR ranks API documentation before every step by BM25 over the instruction and the last two code
blocks, plus a decayed memory of the applications already used (the AppWorld counterpart of the provider
signal $\mathbf{f}_t$), and it completes authentication by adding an application's login API as soon as
that application enters the episode. Committed APIs stay visible. Only the selected documentation enters
the prompt, so the measure of success is whether a small model completes the task within a fixed context
budget. The evaluation is in Chapter~\ref{ch:appworld}; the ToolBench and KCC results above bound the
scope of the claim, because session signals help in genuinely multi-step episodes and act as noise in
queues that switch topic at every step.

%%%%% O2 -> ch5_aprr.tex
\section{Specialisation to AppWorld}
\label{sec:aprr-appworld}

On AppWorld the arms are model tiers rather than agents: a small open model (1.7B--4B parameters) and a
large one (8B--32B), both served on the same GPU. The context records whether the previous step failed,
whether the step is likely to change state, and whether the episode is long, and a Thompson sampler with
a cost penalty on the large tier and a risk margin on the small tier chooses the tier for each step. The
risk margin is the AppWorld form of the risk-signed exploration term: on risky steps the small model must
beat the large one by more than the cost saving before it is trusted. Cost is measured in real tokens;
the evaluation is in Chapter~\ref{ch:appworld}.

%%%%% O3 -> ch6_mncd.tex
\section{Specialisation to AppWorld: Consensus on Outcomes}
\label{sec:mncd-appworld}

When agents act, the object of consensus is not a textual claim but the change an action makes to the
world. On AppWorld, MNCD samples candidate programs for any step that contains a state-changing call,
executes each from a saved world state, records the state-changing calls each one made and whether it
raised an error, restores the state, and certifies the outcome on which a weighted quorum agrees. This
follows the ARC-style practice of generating many programs and keeping the output that independent
programs agree on \citep{chollet2019measure}. The Byzantine reading carries over: a candidate that
raises an error carries no weight, and an outcome is certified only with at least $\min(q,n)$ clean
supporters and a weighted share of at least $\theta$; otherwise MNCD abstains and the router escalates.
The certification threshold and the policy-as-code stage are developed in Chapter~\ref{ch:governance};
the evaluation is in Chapter~\ref{ch:appworld}.

%%%%% O4 -> ch7_fcnp.tex
\section{Specialisation to AppWorld}
\label{sec:fcnp-appworld}

AppWorld episodes run for 30--60 steps, and their outputs (paginated lists, documentation, error
traces) overflow a small model's context. FCNP's typed blocks become the task, recent steps, and older
steps, and its hard-keep set becomes every output line carrying an access token, an identifier or an
error. Older outputs are compressed head-and-tail around those lines, and the PI controller adjusts the
compression length so that the share of critical lines retained stays above the fidelity floor (0.9)
while the history stays within its token budget. The analogue is recency truncation; the evaluation is
in Chapter~\ref{ch:appworld}.
